\documentclass[11pt,a4paper]{article}
% ============================================================
%  Basado en la plantilla de las Prácticas 8-11 (mismos colores/estilos)
% ============================================================

% ---------- Idioma y codificación ----------
\usepackage[utf8]{inputenc}
\usepackage[T1]{fontenc}
\usepackage[spanish,es-noshorthands,es-tabla]{babel}

% ---------- Tipografía y layout ----------
\usepackage{lmodern}
\usepackage[sfdefault,light]{roboto}
\usepackage{geometry}
\geometry{a4paper, top=3.1cm, bottom=2.8cm, left=2.4cm, right=2.4cm, headheight=1.6cm}
\usepackage{ragged2e}
\usepackage{enumitem}
\usepackage{graphicx}
\usepackage{float}
\usepackage{booktabs}
\usepackage{array}
\usepackage[most]{tcolorbox}
\usepackage{tikz}
\usetikzlibrary{shadows,calc}
\usepackage{titlesec}
\usepackage{fancyhdr}
\usepackage{fontawesome5}
\usepackage[colorlinks=true,breaklinks=true]{hyperref}
\usepackage{caption}
\usepackage{microtype}
\usepackage{xurl}
\usepackage{listings}
\usepackage{amsmath,amssymb}
\usepackage{needspace}
\usepackage{icomma}

% ---------- Paleta de colores (tema "coolnight" de tu terminal Kitty) ----------
\definecolor{cNavy}{HTML}{011423}
\definecolor{cTealDeep}{HTML}{064B5A}
\definecolor{cTeal}{HTML}{097289}
\definecolor{cBlue}{HTML}{0FC5ED}
\definecolor{cCyan}{HTML}{0FC5ED}
\definecolor{cMint}{HTML}{44FFB1}
\definecolor{cPurple}{HTML}{A277FF}
\definecolor{cRed}{HTML}{E52E2E}
\definecolor{cBg}{HTML}{EAF6F8}
\definecolor{cText}{HTML}{01111E}
\definecolor{cOutBg}{HTML}{F4F4F4}

\hypersetup{
  linkcolor=cTealDeep,
  citecolor=cTealDeep,
  urlcolor=cBlue,
  pdftitle={Trabajo final: regresión lineal y logística, componentes principales y conglomerados},
  pdfauthor={Miguel Deza}
}

\color{cText}
\emergencystretch=2.5em

% ---------- Encabezado / pie de página ----------
\pagestyle{fancy}
\fancyhf{}
\renewcommand{\headrulewidth}{0pt}
\renewcommand{\footrulewidth}{0pt}
\fancyhead[L]{\small\color{cTealDeep}\textbf{\faGraduationCap\ Maestría en Ciencia de Datos}}
\fancyhead[R]{\small\color{cTeal}Estadística en Ciencia de Datos -- Trabajo final}
\fancyfoot[C]{%
  \begin{tikzpicture}[remember picture]
    \node[circle, fill=cTealDeep, minimum size=1.7em, inner sep=0pt] (pg) {\color{white}\small\thepage};
  \end{tikzpicture}%
}
\renewcommand{\headrule}{{\color{cCyan}\hrule height 1.4pt}\vspace{1pt}{\color{cBlue}\hrule height 0.5pt}}

% ---------- Estilo de secciones ----------
\titleformat{\section}
  {\normalfont\Large\bfseries\color{cTealDeep}}
  {\tikz[baseline=-0.6ex]{\node[fill=cCyan, text=white, rounded corners=1.5pt, inner sep=3pt, font=\bfseries]{\thesection};}}
  {0.6em}{}
  [{\vspace{-0.3em}\color{cCyan}\titlerule[1.2pt]\vspace{0.3em}}]

\titleformat{\subsection}
  {\normalfont\large\bfseries\color{cTeal}}
  {\color{cTeal}\thesubsection}{0.6em}{}

\titlespacing*{\section}{0pt}{1.4em}{0.8em}
\titlespacing*{\subsection}{0pt}{1.1em}{0.5em}

\captionsetup{labelfont={bf,color=cTealDeep}, font=small, textfont=it}

% ---------- Bloques de código R y salida de consola ----------
\lstdefinestyle{styleR}{
  language=R,
  backgroundcolor=\color{cBg},
  basicstyle=\ttfamily\footnotesize\color{cText},
  keywordstyle=\color{cTealDeep}\bfseries,
  commentstyle=\color{gray}\itshape,
  stringstyle=\color{cTeal},
  numbers=none,
  breaklines=true,
  frame=l,
  framerule=2.2pt,
  rulecolor=\color{cCyan},
  xleftmargin=8pt,
  framexleftmargin=6pt,
  aboveskip=6pt, belowskip=6pt,
  showstringspaces=false,
  columns=fullflexible, keepspaces=true
}

\lstdefinestyle{styleOut}{
  backgroundcolor=\color{cOutBg},
  basicstyle=\ttfamily\scriptsize\color{cText},
  breaklines=true,
  frame=l,
  framerule=2.2pt,
  rulecolor=\color{cTealDeep},
  xleftmargin=8pt,
  framexleftmargin=6pt,
  aboveskip=2pt, belowskip=6pt,
  showstringspaces=false,
  columns=fullflexible, keepspaces=true
}

% ---------- Preguntas ----------
\setlist{itemsep=1pt, topsep=3pt}
\graphicspath{{img/}}
\setlength{\parskip}{3pt}

\newtcolorbox{enunciado}[1][]{enhanced, breakable, colback=cBg, colframe=cTealDeep,
  boxrule=0pt, leftrule=3pt, arc=2pt, left=8pt, right=8pt, top=5pt, bottom=5pt,
  fonttitle=\bfseries, coltitle=white, colbacktitle=cTealDeep,
  title={Pregunta #1}}
\newtcolorbox{resultado}{enhanced, breakable, colback=cMint!12, colframe=cTealDeep,
  boxrule=0.5pt, arc=2pt, left=8pt, right=8pt, top=4pt, bottom=4pt}

% \pregunta{etiqueta}{texto del índice}{enunciado}
\newcommand{\pregunta}[3]{%
  \Needspace{12\baselineskip}\medskip
  \phantomsection
  \addcontentsline{toc}{subsection}{#1 \ #2}%
  \begin{enunciado}[#1]#3\end{enunciado}\noindent}

% portada de sección: página completa con fondo oscuro
\newcommand{\seccion}[4]{%
  \clearpage\newgeometry{margin=0cm}
  \refstepcounter{section}\phantomsection
  \addcontentsline{toc}{section}{\protect\numberline{\thesection}\protect\faIcon{#1}\ #2}%
  \thispagestyle{empty}\noindent
  \begin{tikzpicture}[remember picture, overlay]
    \fill[cNavy] (current page.south west) rectangle (current page.north east);
    \fill[cTealDeep, opacity=0.55] (current page.north west) ++(1.5,-1) circle (4.2);
    \fill[cCyan, opacity=0.14] (current page.north west) ++(4.2,-0.4) circle (2.2);
    \fill[cPurple, opacity=0.26] (current page.south east) ++(-1.2,1.4) circle (3.6);
    \fill[cMint, opacity=0.13] (current page.south east) ++(-3.6,0.3) circle (2.0);
    \draw[cCyan, opacity=0.25, line width=0.8pt] (current page.center) ++(5.2,4.6) circle (1.1);
    \node[anchor=center] at ($(current page.center)+(0,5.2)$)
      {\color{cTealDeep!80!cCyan}\bfseries\fontsize{150}{150}\selectfont \thesection};
    \node[circle, fill=cTealDeep, draw=cCyan, line width=1.3pt, minimum size=2.4cm] (ic)
      at ($(current page.center)+(0,1.6)$) {};
    \node at (ic) {\color{cMint}\fontsize{30}{30}\selectfont\faIcon{#1}};
    \node[anchor=center] at ($(current page.center)+(0,-0.4)$)
      {\color{cCyan}\bfseries\fontsize{11}{13}\selectfont S E C C I Ó N};
    \node[anchor=north, text width=14cm, align=center] (tit) at ($(current page.center)+(0,-0.9)$)
      {\hyphenpenalty=10000 {\color{white}\bfseries\fontsize{30}{36}\selectfont #3}};
    \draw[cCyan, line width=1.4pt] ($(tit.south)+(-3.2,-0.3)$) -- ++(6.4,0);
    \node[anchor=north, text width=14cm, align=center] at ($(tit.south)+(0,-0.6)$)
      {\color{cMint}\fontsize{13}{16}\selectfont #4};
  \end{tikzpicture}%
  \clearpage\restoregeometry}

% código y salidas: las líneas de cada bloque se leen de rangos.tex,
% generado por generar-salidas.R desde RLM.R, Logistica.R, CP.R y Conglomerados.R
\expandafter\def\csname r@rlm-datos\endcsname{28-34}
\expandafter\def\csname s@rlm-datos\endcsname{RLM.R}
\expandafter\def\csname r@rlm-ajuste\endcsname{38-41}
\expandafter\def\csname s@rlm-ajuste\endcsname{RLM.R}
\expandafter\def\csname r@rlm-r2\endcsname{45-47}
\expandafter\def\csname s@rlm-r2\endcsname{RLM.R}
\expandafter\def\csname r@rlm-supuestos\endcsname{51-53}
\expandafter\def\csname s@rlm-supuestos\endcsname{RLM.R}
\expandafter\def\csname r@rlm-diag1\endcsname{57-65}
\expandafter\def\csname s@rlm-diag1\endcsname{RLM.R}
\expandafter\def\csname r@rlm-diag2\endcsname{69-77}
\expandafter\def\csname s@rlm-diag2\endcsname{RLM.R}
\expandafter\def\csname r@rlm-vif\endcsname{81-81}
\expandafter\def\csname s@rlm-vif\endcsname{RLM.R}
\expandafter\def\csname r@rlm-backward\endcsname{85-85}
\expandafter\def\csname s@rlm-backward\endcsname{RLM.R}
\expandafter\def\csname r@rlm-final\endcsname{89-93}
\expandafter\def\csname s@rlm-final\endcsname{RLM.R}
\expandafter\def\csname r@rlm-final-sup\endcsname{97-100}
\expandafter\def\csname s@rlm-final-sup\endcsname{RLM.R}
\expandafter\def\csname r@rlm-diag-final\endcsname{104-112}
\expandafter\def\csname s@rlm-diag-final\endcsname{RLM.R}
\expandafter\def\csname r@rlm-sensibilidad\endcsname{116-121}
\expandafter\def\csname s@rlm-sensibilidad\endcsname{RLM.R}
\expandafter\def\csname r@log-datos\endcsname{28-33}
\expandafter\def\csname s@log-datos\endcsname{Logistica.R}
\expandafter\def\csname r@log-justifica\endcsname{37-46}
\expandafter\def\csname s@log-justifica\endcsname{Logistica.R}
\expandafter\def\csname r@log-ajuste\endcsname{50-53}
\expandafter\def\csname s@log-ajuste\endcsname{Logistica.R}
\expandafter\def\csname r@log-or\endcsname{57-58}
\expandafter\def\csname s@log-or\endcsname{Logistica.R}
\expandafter\def\csname r@log-lr\endcsname{62-64}
\expandafter\def\csname s@log-lr\endcsname{Logistica.R}
\expandafter\def\csname r@log-confusion\endcsname{68-74}
\expandafter\def\csname s@log-confusion\endcsname{Logistica.R}
\expandafter\def\csname r@log-roc\endcsname{78-82}
\expandafter\def\csname s@log-roc\endcsname{Logistica.R}
\expandafter\def\csname r@log-roc-fig\endcsname{86-98}
\expandafter\def\csname s@log-roc-fig\endcsname{Logistica.R}
\expandafter\def\csname r@log-cv\endcsname{102-111}
\expandafter\def\csname s@log-cv\endcsname{Logistica.R}
\expandafter\def\csname r@acp-datos\endcsname{28-34}
\expandafter\def\csname s@acp-datos\endcsname{CP.R}
\expandafter\def\csname r@acp-cor\endcsname{38-43}
\expandafter\def\csname s@acp-cor\endcsname{CP.R}
\expandafter\def\csname r@acp-corplot\endcsname{47-51}
\expandafter\def\csname s@acp-corplot\endcsname{CP.R}
\expandafter\def\csname r@acp-pca\endcsname{55-61}
\expandafter\def\csname s@acp-pca\endcsname{CP.R}
\expandafter\def\csname r@acp-scree\endcsname{65-76}
\expandafter\def\csname s@acp-scree\endcsname{CP.R}
\expandafter\def\csname r@acp-cargas\endcsname{80-81}
\expandafter\def\csname s@acp-cargas\endcsname{CP.R}
\expandafter\def\csname r@acp-cargas-fig\endcsname{85-94}
\expandafter\def\csname s@acp-cargas-fig\endcsname{CP.R}
\expandafter\def\csname r@acp-individuos\endcsname{98-114}
\expandafter\def\csname s@acp-individuos\endcsname{CP.R}
\expandafter\def\csname r@acp-biplot\endcsname{118-121}
\expandafter\def\csname s@acp-biplot\endcsname{CP.R}
\expandafter\def\csname r@acp-regresion\endcsname{125-131}
\expandafter\def\csname s@acp-regresion\endcsname{CP.R}
\expandafter\def\csname r@clu-datos\endcsname{28-37}
\expandafter\def\csname s@clu-datos\endcsname{Conglomerados.R}
\expandafter\def\csname r@clu-dist\endcsname{41-44}
\expandafter\def\csname s@clu-dist\endcsname{Conglomerados.R}
\expandafter\def\csname r@clu-dist-fig\endcsname{48-59}
\expandafter\def\csname s@clu-dist-fig\endcsname{Conglomerados.R}
\expandafter\def\csname r@clu-ward\endcsname{63-68}
\expandafter\def\csname s@clu-ward\endcsname{Conglomerados.R}
\expandafter\def\csname r@clu-dendro\endcsname{72-81}
\expandafter\def\csname s@clu-dendro\endcsname{Conglomerados.R}
\expandafter\def\csname r@clu-perfil-hc\endcsname{85-94}
\expandafter\def\csname s@clu-perfil-hc\endcsname{Conglomerados.R}
\expandafter\def\csname r@clu-kmeans\endcsname{98-109}
\expandafter\def\csname s@clu-kmeans\endcsname{Conglomerados.R}
\expandafter\def\csname r@clu-centroides\endcsname{113-121}
\expandafter\def\csname s@clu-centroides\endcsname{Conglomerados.R}
\expandafter\def\csname r@clu-compara\endcsname{125-136}
\expandafter\def\csname s@clu-compara\endcsname{Conglomerados.R}
\expandafter\def\csname r@clu-plano\endcsname{140-151}
\expandafter\def\csname s@clu-plano\endcsname{Conglomerados.R}
\expandafter\def\csname r@clu-enfermedad\endcsname{155-161}
\expandafter\def\csname s@clu-enfermedad\endcsname{Conglomerados.R}
\expandafter\def\csname r@clu-enf-fig\endcsname{165-173}
\expandafter\def\csname s@clu-enf-fig\endcsname{Conglomerados.R}

\newcommand{\codigo}[1]{%
  \edef\figopts{linerange={\csname r@#1\endcsname}}\edef\figfile{\csname s@#1\endcsname}%
  \expandafter\lstinputlisting\expandafter[\figopts, style=styleR]{\figfile}}
\newcommand{\salida}[1]{\lstinputlisting[style=styleOut]{salidas/#1.txt}}
% figura a la izquierda y su código R a la derecha
\newcommand{\figcod}[2]{%
  \par\medskip\noindent
  \begin{minipage}[c]{0.47\linewidth}\centering
    \includegraphics[width=\linewidth]{#1}\captionof{figure}{#2}
  \end{minipage}\hfill
  \begin{minipage}[c]{0.5\linewidth}
    \edef\figopts{linerange={\csname r@#1\endcsname}}\edef\figfile{\csname s@#1\endcsname}%
    \expandafter\lstinputlisting\expandafter[\figopts, style=styleR, basicstyle=\ttfamily\scriptsize\color{cText}]{\figfile}
  \end{minipage}\par\medskip}
\newcommand{\res}[1]{\begin{resultado}\textbf{\color{cTealDeep}Resultado.} #1\end{resultado}}
\newcommand{\var}[1]{\texttt{\detokenize{#1}}}

% ============================================================
% portada: \portada{icono}{título}{subtítulo}{descripción}
\newcommand{\portada}[4]{%
\begin{titlepage}
\newgeometry{margin=0cm}
\begin{tikzpicture}[remember picture, overlay]
  \fill[cNavy] (current page.south west) rectangle (current page.north east);
  \fill[cPurple, opacity=0.28] (current page.north east) ++(-1,1) circle (3.2);
  \fill[cCyan, opacity=0.22] (current page.north east) ++(-2.6,0.4) circle (2.1);
  \fill[cTealDeep, opacity=0.55] (current page.south west) ++(1,-1) circle (4.0);
  \fill[cMint, opacity=0.16] (current page.south west) ++(2.4,-0.2) circle (2.3);

  \node[circle, fill=cTealDeep, draw=cCyan, line width=1.3pt, minimum size=2.5cm] (logo) at ($(current page.north)+(0,-2.9)$) {};
  \node at (logo) {\color{cMint}\fontsize{28}{28}\selectfont\faIcon{#1}};

  \node[anchor=north, text width=14cm, align=center] at ($(current page.north)+(0,-4.7)$) {
    {\color{cCyan}\bfseries\fontsize{12}{14}\selectfont \faDatabase\ \ MAESTRÍA EN CIENCIA DE DATOS}
  };

  \node[anchor=center, text width=15cm, align=center] at ($(current page.center)+(0,1.2)$) {
    {\color{white}\bfseries\fontsize{30}{38}\selectfont #2}\\[0.35em]
    {\color{cMint}\bfseries\fontsize{17}{21}\selectfont \faIcon{#1}\ #3}
  };

  \draw[cCyan, line width=1.4pt] ($(current page.center)+(-3.6,-0.9)$) -- ++(7.2,0);

  \node[anchor=center, text width=13cm, align=center] at ($(current page.center)+(0,-1.9)$) {
    {\color{white!85!black}\fontsize{11}{14}\selectfont
     #4}
  };

  \node[anchor=south, text width=14cm, align=center] at ($(current page.south)+(0,2.6)$) {
    \begin{tabular}{r l}
      \color{cCyan}\bfseries \faUser\ Elaborado por: & \color{white} Miguel Deza \\[3pt]
      \color{cCyan}\bfseries \faGraduationCap\ Programa: & \color{white} Maestría en Ciencia de Datos -- UNSA \\[3pt]
      \color{cCyan}\bfseries \faBook\ Curso: & \color{white} Estadística en Ciencia de Datos \\[3pt]
      \color{cCyan}\bfseries \faChalkboardTeacher\ Docente: & \color{white} Carlos Antonio Abanto Valle \\[3pt]
      \color{cCyan}\bfseries \faIcon{calendar-alt}\ Fecha: & \color{white} 7 de noviembre de 2026 \\
    \end{tabular}
  };
\end{tikzpicture}
\end{titlepage}
\restoregeometry
}

\hypersetup{pdftitle={Análisis de componentes principales}}
\begin{document}
\portada{project-diagram}{Análisis de\\ Componentes Principales}{Reducción de dimensionalidad con la matriz de correlaciones}{Base de datos: Student Performance (UCI)\\ Correlaciones, scree plot, cargas y regresión sobre componentes}

\setcounter{tocdepth}{2}
\tableofcontents

\vspace{1.2em}
\noindent Todo el análisis está en \var{3_Componentes_Principales.R} (R 4.6.1; paquetes \var{corrplot} y \var{car}), que lee \var{student.csv} (395 estudiantes). En cada pregunta muestro el bloque de código, su salida de consola y la interpretación.

\clearpage
\section*{Variables}
Uso las 11 variables cuantitativas del enunciado, todas como numéricas:
\codigo{acp-datos}
\salida{acp-datos}

\pregunta{a)}{Por qué estandarizar}{Explique por qué es necesario estandarizar las variables antes de aplicar el
ACP en este contexto.}
El ACP busca las direcciones de máxima varianza, así que una variable con varianza grande domina las
primeras componentes solo por su escala. Aquí \var{inasistencias} va de 0 a 75 con $s=8,0$ ($s^2=64$),
mientras que las escalas Likert 1--5 tienen $s\approx1$ y \var{reprobaciones_previas} $s=0,74$ ($s^2=0,55$).
Sin estandarizar, la CP1 sería prácticamente ``inasistencias'' (64 de unas 76 unidades de varianza
total, el 84\,\%). Además las unidades no son comparables (años, faltas, escalas ordinales). Con
\[
z_{ij}=\frac{x_{ij}-\bar x_j}{s_j}
\]
todas tienen varianza 1 y el ACP sobre $Z$ equivale a diagonalizar la matriz de correlaciones $R$.

\res{Estandarizo para que cada variable pese lo mismo y las componentes reflejen correlaciones, no unidades.}

\pregunta{b)}{Matriz de correlaciones}{Calcule la matriz de correlaciones entre las variables y comente
brevemente la presencia de relaciones lineales relevantes.}
\codigo{acp-cor}
\salida{acp-cor}
\figcod{acp-corplot}{Matriz de correlaciones (triángulo superior).}
Relaciones más relevantes:
\begin{itemize}[leftmargin=1.2em]
  \item[\color{cTeal}\faLink] \textbf{alcohol días de semana -- fin de semana}: $r=0,65$, la más fuerte.
  \item[\color{cTeal}\faLink] \textbf{educación de la madre -- del padre}: $r=0,62$.
  \item[\color{cTeal}\faLink] \textbf{salidas sociales} con alcohol fin de semana ($0,42$), tiempo libre ($0,29$)
  y alcohol semana ($0,27$): un bloque de ``vida social''.
  \item[\color{cRed}\faUnlink] \textbf{reprobaciones} con educación de los padres ($-0,24$, $-0,25$) y con
  edad ($+0,24$); \textbf{tiempo de estudio} con alcohol fin de semana ($-0,25$).
\end{itemize}
El resto está por debajo de $|0,2|$ y \var{salud} casi no se relaciona con nada. El test de esfericidad de
Bartlett ($\chi^2_{55}=701,3$, $p\approx10^{-113}$) rechaza $R=I$: hay correlación suficiente para que el ACP
tenga sentido, aunque las correlaciones moderadas anticipan que la varianza se repartirá en varias componentes.

\pregunta{c)}{ACP con la matriz de correlaciones}{Aplique el Análisis de Componentes Principales utilizando la
matriz de correlaciones.}
\[
R\,\mathbf a_k=\lambda_k\,\mathbf a_k,\qquad \sum_{k=1}^{11}\lambda_k=\operatorname{tr}(R)=11,\qquad
\text{CP}_k=Z\,\mathbf a_k
\]
\codigo{acp-pca}
\salida{acp-pca}
Con \var{prcomp(scale. = TRUE)} los autovalores son los de $R$ (suman 11) y los vectores \var{rotation}
son las cargas $\mathbf a_k$.

\pregunta{d)}{Scree plot y número de componentes}{Presente el scree plot e indique cuántos componentes deberían
retenerse según la varianza explicada acumulada (80\,\% o 90\,\%).}
\figcod{acp-scree}{Scree plot (arriba) y varianza acumulada con las líneas de 80\,\% y 90\,\% (abajo).}
\begin{center}\small
\begin{tabular}{lcc}
\toprule
Criterio & Componentes & Varianza retenida \\ \midrule
Varianza acumulada $\ge$ 80\,\% & \textbf{7} & 81,9\,\% \\
Varianza acumulada $\ge$ 90\,\% & 9 & 94,0\,\% \\
Kaiser ($\lambda_k>1$) & 4 & 58,3\,\% \\
Codo del scree plot & 3 & 48,8\,\% \\ \bottomrule
\end{tabular}
\end{center}
Retengo \textbf{7 componentes} (criterio del 80\,\%): con 9 llegaría al 94\,\%, pero pasaría de 11 a 9
variables, que casi no reduce la dimensión. La curva cae despacio después de CP2, consecuencia de las
correlaciones bajas: este conjunto no se comprime bien en pocas componentes.

\res{7 CP según el 80\,\% (81,9\,\%); 9 según el 90\,\%.}

\pregunta{e)}{Interpretación de CP1 y CP2}{Interprete los dos primeros componentes en términos de las variables
originales, analizando los signos y magnitudes de las cargas.}
\codigo{acp-cargas}
\salida{acp-cargas}
\figcod{acp-cargas}{Cargas de las dos primeras componentes.}
\begin{itemize}[leftmargin=1.2em]
  \item[\color{cTeal}\faGlassCheers] \textbf{CP1 (21,1\,\%) -- ``estilo de vida social / riesgo''}: cargas
  negativas grandes en alcohol fin de semana ($-0,52$), alcohol semana ($-0,48$), salidas sociales ($-0,37$),
  reprobaciones ($-0,30$), tiempo libre ($-0,27$) y edad ($-0,24$), y positiva en tiempo de estudio ($+0,28$).
  Un CP1 \emph{bajo} es un alumno que bebe, sale mucho, estudia poco, es mayor y ha reprobado; un CP1
  \emph{alto} es el perfil estudioso. (El signo global es arbitrario.)
  \item[\color{cTeal}\faUserGraduate] \textbf{CP2 (16,6\,\%) -- ``capital educativo familiar''}: dominada por
  educación de la madre ($+0,61$) y del padre ($+0,60$), contra reprobaciones ($-0,30$) y edad ($-0,22$). Un
  CP2 alto es un alumno de padres con educación superior y sin repitencia.
\end{itemize}
\var{salud} e \var{inasistencias} casi no cargan en ninguna de las dos; aparecen en CP3 (inasistencias
$-0,61$, salud $+0,44$).

\pregunta{f)}{Individuos en el plano CP1--CP2}{Represente gráficamente los individuos en el plano definido por
los dos primeros componentes e identifique posibles patrones o agrupamientos.}
\figcod{acp-individuos}{Puntajes de los 395 alumnos en CP1--CP2.}
\figcod{acp-biplot}{Biplot: individuos (gris) y cargas (flechas).}
\begin{itemize}[leftmargin=1.2em]
  \item[\color{cTeal}\faEye] No hay grupos separados: la nube es continua y asimétrica, con una cola hacia
  CP1 negativo.
  \item[\color{cRed}\faEye] Esa cola son los alumnos con alto consumo de alcohol el fin de semana (rojo): su
  CP1 medio es $-1,98$ frente a $0,50$ del resto. El biplot muestra las flechas de alcohol y salidas
  apuntando justo hacia ella.
  \item[\color{cTeal}\faEye] Por aprobación los puntos se mezclan, pero los aprobados están algo desplazados
  hacia arriba y a la derecha (CP1 medio $0,23$ vs.\ $-0,47$; CP2 $0,15$ vs.\ $-0,30$): estudiar más, beber
  menos y tener padres más educados va con aprobar.
\end{itemize}

\pregunta{g)}{CP como variables explicativas}{Discuta si los componentes principales obtenidos podrían ser
utilizados posteriormente como variables explicativas en un modelo de regresión.}
Sí: es la \emph{regresión sobre componentes principales}. Lo probé con las 7 CP retenidas para explicar la nota
final:
\codigo{acp-regresion}
\salida{acp-regresion}
\begin{itemize}[leftmargin=1.2em]
  \item[\color{cTeal}\faPlusCircle] \textbf{Ventaja}: las CP son ortogonales (VIF $=1$ exacto), así que
  eliminan cualquier multicolinealidad y cada coeficiente se estima de forma independiente.
  \item[\color{cTeal}\faPlusCircle] CP1 y CP2 resultan significativas y con signo esperado: más ``estilo
  estudioso'' y más capital educativo, mejor nota.
  \item[\color{cRed}\faMinusCircle] \textbf{Desventaja 1}: las CP se eligen por varianza de las $X$, no por
  su relación con $Y$. Aquí CP5 y CP7, de poca varianza, son significativas, y CP3--CP4 no aportan.
  \item[\color{cRed}\faMinusCircle] \textbf{Desventaja 2}: interpretar ``+1 en CP7'' es difícil, y con 7 de
  11 dimensiones se pierde algo de ajuste ($R^2_{aj}$: 0,137 frente a 0,153 con las 11 variables originales).
  \item[\color{cRed}\faMinusCircle] En esta base la multicolinealidad no era un problema (en el ejercicio de regresión lineal
  sobre esta base los VIF fueron $<1,8$), por lo que el beneficio principal del ACP no aplica.
\end{itemize}

\res{Se pueden usar y garantizan ortogonalidad, pero aquí convienen poco: las correlaciones son bajas, se pierde
interpretabilidad y ajuste. Las CP serían útiles con muchas variables muy correlacionadas o con $p$ cercano a
$n$; una alternativa es PLS, que elige componentes pensando en $Y$.}

\end{document}
